\begin{table}[htbp]
\centering
\caption{Counterfactual female shares among high-achieving students assigned to STEM}
\label{tab:counterfactuals}
\footnotesize
\setlength{\tabcolsep}{4pt}
\begin{tabular}{lcccc}
\toprule
 & & \multicolumn{2}{c}{Counterfactual} & Change from \\
\cmidrule(lr){3-4}
Setting & Observed & Balanced pipeline & \shortstack{Common STEM\\assignment rates} & balancing \\
\midrule
Uganda & 26.3 & 36.2 & 38.6 & +9.9 \\
Chile & 27.5 & 32.5 & 44.0 & +5.1 \\
Taiwan & 29.3 & 38.8 & 39.5 & +9.5 \\
Brazil & 31.5 & 32.0 & 49.4 & +0.5 \\
China & 33.3 & 28.8 & 55.3 & $-$4.6 \\
Finland & 35.1 & 33.6 & 51.6 & $-$1.5 \\
Sweden & 42.2 & 27.5 & 65.9 & $-$14.7 \\
Australia & 42.8 & 39.2 & 53.7 & $-$3.6 \\
Spain & 44.8 & 32.5 & 62.8 & $-$12.3 \\
Greece & 45.1 & 34.4 & 61.1 & $-$10.7 \\
\bottomrule
\end{tabular}
\begin{minipage}{\textwidth}\vspace{4pt}\scriptsize \textit{Notes:} All entries are the percentage of high-achieving students assigned to a STEM program who are women, where high-achieving means the top decile of each setting's academic performance distribution. \textit{Observed} is the realized share. \textit{Balanced pipeline} holds each gender's rate of reaching a STEM assignment at its observed value and sets the top decile to equal numbers of women and men; it equals $1-\Lambda(C+V)$ in the notation of \eqref{eq:decomposition_assign} (Online \ref{sec:appendix_decomposition}). \textit{Common STEM assignment rates} holds the observed top-decile composition and gives women men's assignment rate $M_m/T_m$, setting $C+V=0$; the female share is then $1-\Lambda(P)$. The final column is the balanced-pipeline share minus the observed share; a negative entry means that balancing the pipeline, on its own, would lower the female share among STEM assignees, because women already outnumber men in that setting's top decile. Both columns are accounting exercises that reallocate observed rates and are not predictions of how either margin would respond to an intervention.\end{minipage}
\end{table}
